\section{方法比较与适用场景分析}\label{sec:comparison}

\subsection{统一维度下的横向比较}\label{subsec:unified-comparison}

横向比较需区分直接改变的变量与最终系统指标。图~\ref{fig:communication-cost-map}把式~\eqref{eq:total-bits}中的轮数、参与集合和上传载荷分别对应频率控制、客户端选择和压缩；速率与时延、并发聚合及消息路径则对应资源分配、AirComp和分层联邦学习。

\begin{figure*}[t]
  \centering
  \resizebox{0.96\textwidth}{!}{\begin{tikzpicture}[
  font=\small,
  >=Latex,
  root/.style={draw, rounded corners=2pt, minimum width=5.2cm,
               minimum height=0.78cm, align=center, fill=black!9,
               font=\bfseries},
  item/.style={draw, rounded corners=1.5pt, text width=2.25cm,
               minimum height=1.15cm, align=center, fill=black!3,
               inner sep=3pt},
  flow/.style={->, line width=0.55pt}
]
  \node[root] (root) at (0,1.75) {联邦训练的通信成本与完成代价};
  \coordinate (bus) at (0,0.95);
  \draw[line width=0.55pt] (root.south) -- (bus);
  \draw[line width=0.55pt] (-6.45,0.95) -- (6.45,0.95);

  \node[item] (round) at (-6.45,0) {全局轮数 \(T\)\\[2pt]\textbf{通信频率控制}};
  \node[item] (client) at (-3.87,0) {参与集合 \(\mathcal S_t\)\\[2pt]\textbf{客户端选择}};
  \node[item] (payload) at (-1.29,0) {单次载荷 \(B_{k,t}^{\mathrm{ul}}\)\\[2pt]\textbf{量化与稀疏化}};
  \node[item] (radio) at (1.29,0) {速率 \(R\)、时延、能耗\\[2pt]\textbf{无线资源分配}};
  \node[item] (aircomp) at (3.87,0) {同时接入与聚合误差\\[2pt]\textbf{AirComp}};
  \node[item] (hier) at (6.45,0) {消息路径与回传次数\\[2pt]\textbf{分层联邦学习}};

  \foreach \n in {round,client,payload,radio,aircomp,hier}{
    \draw[flow] (\n |- bus) -- (\n.north);
  }
\end{tikzpicture}
}
  \caption{通信成本分量与五类优化方法的对应关系\protect\cite{Lim2020MobileEdge,Zhang2023EdgeFLSurvey,Ge2025EfficientFLSurvey}}
  \label{fig:communication-cost-map}
\end{figure*}

同一方法可能影响多个指标，但分类应以其直接控制的变量为准：压缩改变比特，资源分配改变速率和等待。表~\ref{tab:method-comparison}据此按“变量—原理—收益—代价”比较五类方法。

\begin{table*}[t]
  \centering
  \caption{五类通信效率优化方法的统一比较}
  \label{tab:method-comparison}
  \begin{threeparttable}
    \fontsize{7.1pt}{8.6pt}\selectfont
    \setlength{\tabcolsep}{2.8pt}
    \renewcommand{\arraystretch}{1.22}
    \begin{tabularx}{\textwidth}{
      @{}>{\centering\arraybackslash}p{1.35cm}
      >{\centering\arraybackslash}p{1.25cm}
      >{\raggedright\arraybackslash}p{1.72cm}
      >{\raggedright\arraybackslash}p{1.62cm}
      >{\raggedright\arraybackslash}p{1.62cm}
      >{\raggedright\arraybackslash}p{1.68cm}
      >{\raggedright\arraybackslash}p{1.62cm}
      >{\raggedright\arraybackslash}X@{}}
      \toprule
      方法类别 & 优化对象 & 代表机制 & 主要收益 & 额外代价 & 关键假设 & 适用场景 & 主要局限 \\
      \midrule
      客户端选择与频率控制
        & \(T,\mathcal S_t,E_k\)
        & 效用选择；异步或事件触发
        & 少上传、少同步、少等待
        & 状态上报、偏差与陈旧
        & 状态可观测；更新具代表性
        & 大规模异构物联网
        & 弱设备可能长期缺席，影响公平与收敛\cite{Nishio2019FedCS,Wang2019AdaptiveFL,Zhang2024ClientSelectionSurvey} \\
      无线资源分配
        & 带宽、功率、时隙、CPU频率及边缘关联
        & 跨轮资源与关联联合调度
        & 降低时延、能耗和掉队率
        & 信道反馈、求解和信令
        & 状态可估计；链路模型有效
        & 蜂窝、多业务、无人机网络
        & 改变传输效率但未必减少累计比特\cite{Dinh2021WirelessConvergence,Xu2022MultiFLBandwidth,Fan2025EdgeAssociation,Meng2025UAVEnergy} \\
      模型与梯度压缩
        & \(B_{k,t}^{\mathrm{ul}}\)、稀疏度和量化位宽
        & 量化、稀疏、误差反馈
        & 减少载荷和发送能耗
        & 失真、索引和重构计算
        & 更新可压缩；误差可补偿
        & 高维模型、窄带终端
        & 理论压缩率不等于有效传输比\cite{Reisizadeh2020FedPAQ,Sattler2020STC,Marnissi2024AdaptiveSQ} \\
      AirComp
        & 多用户上传时间、频谱及聚合误差
        & 同步模拟发送；波形域求和
        & 传输与聚合并行
        & 信道状态、同步、功控、噪声
        & 函数可加；物理层可协同
        & 密集同步无线接入
        & 深衰落和时变信道会破坏对齐\cite{Zhu2020BroadbandAnalog,Tegin2023TimeVaryingOTA,Oh2024SparseOneBitOTA} \\
      分层联邦学习
        & 消息路径、簇划分及两级聚合间隔 \(H\)
        & 边缘高频、云端低频聚合
        & 少回传、易扩展
        & 部署、分组、控制和漂移
        & 边缘可用；簇权重合理
        & 多小区、云—边—端集群
        & 簇偏差、节点迁移和故障增加协调难度\cite{Liu2020HierarchicalFL,Liu2023HierarchicalQuantization,Wang2023HierarchicalIoT,Zhong2022HierarchicalHeterogeneous} \\
      \bottomrule
    \end{tabularx}
    \begin{tablenotes}[flushleft]
      \footnotesize
      \item 注：\(T\) 为全局轮数，\(\mathcal S_t\) 为参与集合，\(E_k\) 为本地步数；表中内容不用于比较不同实验设置中的百分比。
    \end{tablenotes}
  \end{threeparttable}
\end{table*}


表~\ref{tab:method-comparison}表明，通信收益通常伴随成本转移：少参与带来偏差风险，压缩和AirComp引入失真，资源与分层优化增加控制开销。比特减少不必然缩短端到端时延，单轮加速也不保证总能耗下降\cite{Tran2019WirelessOptimization,Ge2025EfficientFLSurvey}，因而不能脱离目标精度和无线基线判断最优类别。

\subsection{场景匹配与方法组合}\label{subsec:scenario-matching}

方法选择应先识别主导瓶颈，再决定是否组合多个优化层。表~\ref{tab:scenario-selection}中的“优先”只表示先处理主导成本。若设备可用性还受参与意愿影响，激励应作为客户端选择的约束，而非新的通信类别\cite{Jia2023HierarchicalIncentive}。

\begin{table*}[t]
  \centering
  \caption{不同无线边缘场景下的方法选择}
  \label{tab:scenario-selection}
  \begin{threeparttable}
    \small
    \setlength{\tabcolsep}{4pt}
    \renewcommand{\arraystretch}{1.28}
    \begin{tabularx}{\textwidth}{
      @{}>{\centering\arraybackslash}p{2.05cm}
      >{\raggedright\arraybackslash}p{2.45cm}
      >{\centering\arraybackslash}p{2.65cm}
      >{\raggedright\arraybackslash}X
      >{\raggedright\arraybackslash}X@{}}
      \toprule
      主导瓶颈 & 典型场景 & 优先考虑的方法 & 选择理由 & 使用时需要注意的问题 \\
      \midrule
      同步等待与终端异构
        & 大规模物联网；设备能力和在线状态差异明显
        & 客户端选择＋资源分配
        & 同时控制“谁参与”和“如何完成本轮传输”，缓解慢终端造成的等待
        & 快设备优先会造成选择偏差；状态采集和联合求解也有开销\cite{Xu2021ClientBandwidth,Nishio2019FedCS} \\
      高维更新与窄上行
        & 视觉模型、移动终端或低速上行链路
        & 量化、稀疏化及误差反馈
        & 直接减少每轮需要编码和发送的更新量
        & 需计入索引、码本和残差；压缩误差可能增加训练轮数\cite{Reisizadeh2020FedPAQ,Sattler2020STC} \\
      密集终端同时接入
        & 同步条件可控、聚合函数可加的蜂窝或传感网络
        & AirComp；必要时结合稀疏化
        & 在同一时频资源中完成多用户叠加，避免逐用户正交接入
        & 依赖信道估计、同步与功率可行性，并引入聚合误差\cite{Zhu2020BroadbandAnalog,Oh2024SparseOneBitOTA} \\
      回传链路受限
        & 多小区、工业园区或云—边—端集群
        & 分层联邦＋两级压缩
        & 先减少边—云同步频率，再降低簇模型的回传载荷
        & 需联合控制量化误差、簇间漂移和边缘负载\cite{Liu2023HierarchicalQuantization,Zhong2022HierarchicalHeterogeneous} \\
      多瓶颈随时间变化
        & 移动、多业务竞争或时变信道环境
        & 跨轮选择、压缩和资源联合优化
        & 可随训练阶段和网络状态切换主要优化变量
        & 在线决策、反馈和模型迁移可能抵消部分通信收益\cite{Xu2021ClientBandwidth,Marnissi2024AdaptiveSQ,Azimi2024HierarchicalOTA} \\
      \bottomrule
    \end{tabularx}
  \end{threeparttable}
\end{table*}


客户端选择与无线资源分配的组合同时回答“选谁”和“如何传”。前者筛除本轮难以完成训练的终端，后者为入选设备配置带宽、功率和时隙；但两层若都偏向快设备，会放大数据代表性偏差，状态上报与在线求解也会形成控制流量\cite{Xu2021ClientBandwidth,Nishio2019FedCS}。

稀疏化与AirComp同时降低有效维度和多用户接入时间。终端需协调稀疏位置或采用可聚合编码，否则不同支持集无法按坐标叠加；量化误差、残差近似、信道失配和噪声还会共同进入聚合误差\cite{Oh2024SparseOneBitOTA}。

分层联邦与压缩先以低频云端同步减少回传次数，再压缩簇模型以降低单次回传载荷。两级压缩误差会与簇间漂移叠加，边缘侧还要承担编解码、簇管理和负载均衡\cite{Liu2023HierarchicalQuantization,Chen2025EdgeAssistedSparsification}。因此，应分别核算端—边和边—云开销，不能直接相加不同基线下的收益。

\section{现存问题与发展趋势}\label{sec:challenges}

\subsection{现有研究的共同局限}\label{subsec:common-limitations}

指标口径不统一会限制横向比较。轮数、压缩率、时延、能耗和聚合误差衡量的对象不同；若目标精度以及控制消息、索引和重传的统计范围不同，相同“通信降低”并不表示相同系统收益\cite{Lim2020MobileEdge,Mills2020WirelessEdgeIoT,Ge2025EfficientFLSurvey}。

理想信道假设和终端异构都会限制部署结论。固定拓扑、准确信道状态和同步终端便于分析，但移动、深衰落和故障会改变参与集合、AirComp对齐与分层归属\cite{Tegin2023TimeVaryingOTA,Azimi2024HierarchicalOTA}。设备能力还可能与数据分布相关，快设备优先、长期本地训练和固定分簇都可能造成统计偏差。

跨层组合会叠加误差和控制开销：压缩失真、信道噪声、更新陈旧与簇间漂移可能同时出现，状态反馈和在线求解也会消耗资源。区块链异步方案还需执行模型上链、共识和信誉评估，产生额外通信与计算\cite{Huang2024BlockchainAFL}。秘密共享和安全传输同样只作为效率约束讨论，不构成独立主线\cite{Chen2022SecureCompressedFL,Meng2025UAVEnergy}。

\subsection{面向可部署系统的发展方向}\label{subsec:future-directions}

发展方向应回应上述限制。在线跨层方法需要让客户端选择、压缩程度、无线资源和边缘关联随训练阶段与信道变化，同时计入决策时间和反馈比特\cite{Xu2021ClientBandwidth,Marnissi2024AdaptiveSQ,Azimi2024HierarchicalOTA}。面向异构终端的方法还要同时处理设备能力、模型规模和数据分布差异。边缘辅助与集群蒸馏依赖额外边缘算力或代理知识，其收益仍需跨场景验证\cite{Zhong2022HierarchicalHeterogeneous,Chen2025EdgeAssistedSparsification,Zhang2026FedCD}。

面向任务的消息设计不再要求每轮传输完整参数。单轮通信方案让客户端上传生成模型，由服务器生成模拟数据并训练全局模型\cite{Zheng2023SingleRound}；集群蒸馏则交换任务知识\cite{Zhang2026FedCD}。这类证据仍受具体任务和生成质量限制。可复现实验应在相同目标精度下报告累计比特、墙钟时延、能耗、控制开销和聚合误差，并公开信道轨迹、数据划分与重复设置。

\section{结论}\label{sec:conclusion}

通信效率方法可归纳为客户端选择与频率控制、无线资源分配、模型压缩、AirComp和分层拓扑五类。它们分别作用于轮数与参与节点、传输载荷、无线过程和消息路径，也会引入统计偏差、失真、控制信令或部署复杂度。窄带高维模型可优先考虑压缩；密集且同步条件可控的接入场景可考虑AirComp；多小区或回传受限系统更适合分层聚合。组合方法应先确定主导瓶颈，再在同一目标精度下分别核算比特、时延、能耗和误差，不能把不同实验基线下的收益直接相加。
